\chapter{液体的性质}
\problem{6.1}{研究表明，液体的黏度随温度变化的行为可以表示为
\[
\eta=\eta_0e^{E_a/(k_{\rm B}T)}\qquad\text{（6.4 式）},
\]
其中 $\eta_0$ 为常量，$E_a$ 为液体的组分粒子的激活能，$k_{\rm B}$ 为玻尔兹曼常量，$T$ 为热力学温标下的温度。现有实验测得，在其冰点，水的黏度为 $1.793\times10^{-3}\,\mathrm{N\cdot s/m^2}$，在其汽点，水的黏度为 $2.84\times10^{-4}\,\mathrm{N\cdot s/m^2}$。试近似确定上式中的常量。}
\answerarea
\problem{6.2}{对液体作范德瓦尔斯模型近似，试导出液体的摩尔熵随系统的温度和体积变化的行为，并给出液体的黏度与摩尔熵的比值的表达式，画出其随温度变化的行为的曲线。}
\answerarea
\clearpage
\problem{6.3}{一实验中测得水的自扩散系数为 $1.5\times10^{-9}\,\mathrm{m^2/s}$，试由之估算水分子的定居时间。}
\answerarea
\problem{6.4}{在 $20\,\mathrm{km^2}$ 的湖面上，下了一场降雨量为 $50\,\mathrm{mm}$ 的大雨，雨滴半径 $r=2.0\,\mathrm{mm}$。设温度不变，已知水的表面张力系数 $\sigma=7.3\times10^{-2}\,\mathrm{N/m}$，该下雨过程中释放出了多少能量？}
\answerarea
\clearpage
\problem{6.5}{实验测得水的表面张力系数为 $0.073\,\mathrm{N/m}$，如果将两个半径都为 $2\,\mathrm{mm}$ 的水滴合并为一个大水滴，则水滴的温度最多可以改变多少？}
\answerarea
\problem{6.6}{将一根表面涂有油脂的钢针轻轻地放到水面上，如果水的表面张力系数为 $0.073\,\mathrm{N/m}$，钢的密度为 $7.8\times10^3\,\mathrm{kg/m^3}$，要保证钢针不下沉，钢针的最大直径为多大？}
\answerarea
\clearpage
\problem{6.7}{试利用卡诺定理证明确定液体表面内能的（6.14）式：
\begin{equation}
U_s=\sigma A_s-TA_s\left(\frac{\partial\sigma}{\partial T}\right)_{A_s}.
\tag{6.14}
\end{equation}}
\answerarea
\problem{6.8}{设有一个在 $1.0136\times10^5\,\mathrm{Pa}$ 大气中吹成的、半径为 $1.0\times10^{-2}\,\mathrm m$ 的球形肥皂泡。已知肥皂膜的表面张力系数为 $\sigma=5.0\times10^{-2}\,\mathrm{N/m}$，试确定周围的大气压强多大，才可使泡的半径在等温条件下增大为 $2.0\times10^{-2}\,\mathrm m$。}
\answerarea
\clearpage
\problem{6.9}{两个表面张力系数都为 $\sigma$、半径分别为 $R_1,R_2$ 的肥皂泡，在相同的大气中，等温地并合成一个半径为 $R$ 的肥皂泡，试证明：如果泡内的气体可以近似为理想气体，则泡外大气的压强为
\[
p=\frac{4\sigma(R^2-R_1^2-R_2^2)}{R_1^3+R_2^3-R^3}.
\]}
\answerarea
\problem{6.10}{测定液体表面张力系数的一种简便方法是称量从毛细管下端滴出液体的重量，并利用快速连续摄影方法，测定液滴在脱离毛细管下端的瞬间液滴“颈”的直径 $d$。实验测定 $d=0.7\,\mathrm{mm}$，318 滴液滴的总重量为 $5\,\mathrm g$。试问此种液体的表面张力系数是多少？}
\answerarea
\clearpage
\problem{6.11}{我们经常说水面上泛起漂亮的涟漪，所谓的涟漪就是水面波，水面波也常被称为表面张力波，因为它起因于表面张力。如果实验观测记录下水面波的频率、波长和液体密度分别为 $\nu,\lambda$ 和 $\rho$，试确定水的表面张力系数。}
\answerarea
\problem{6.12}{利用检验弯曲液面内外压强差的拉普拉斯公式的图 6.15 所示实验装置，在其两端吹肥皂泡。
\begin{enumerate}
\item 试证明：在肥皂泡的形状小于半球的情况下，吹成的肥皂泡是稳定的；如果吹成的肥皂泡的形状大于半球，则肥皂泡不稳定。
\item 如果装置中两管的管口直径都为 $4\,\mathrm{mm}$，吹气停止时形成的肥皂泡的直径都为 $5\,\mathrm{mm}$，假设因肥皂膜的位置移动引起的泡内空气的密度的改变可以忽略不计，试确定稳定的肥皂泡的状态。
\end{enumerate}
\bookfig{6_12}{图6.15\quad 检验拉普拉斯公式的实验装置}}
\answerarea
\clearpage
\problem{6.13}{一直径为 $0.1\,\mathrm{mm}$ 的球形气泡恰好处在水面下，如果水面上的大气压强为 $1.0\times10^5\,\mathrm{Pa}$，试确定气泡内的压强。}
\answerarea
\problem{6.14}{记液体的等温压缩系数为 $\kappa$、液体与空气之间的表面张力系数为 $\sigma$，
\begin{enumerate}
\item 试确定半径为 $r$ 的液滴的密度。
\item 实验测量结果表明，在温度为 $4\,^\circ\mathrm C$ 的情况下，水的等温压缩系数 $\kappa=4.75\times10^{-7}\,\mathrm{atm^{-1}}$，水与空气之间的表面张力系数为 $\sigma=7.24\times10^{-2}\,\mathrm{N\cdot m^{-1}}$，试确定 $4\,^\circ\mathrm C$ 情况下，半径为 $100\,\mu\mathrm m$ 的小水滴的重量。
\end{enumerate}}
\answerarea
\clearpage
\problem{6.15}{在压强为 $p_0$ 的大气中有一半径为 $R_0$ 的电中性肥皂泡，记肥皂液的表面张力系数为 $\sigma$，
\begin{enumerate}
\item 这种情况下，肥皂泡中的气体的压强为多大？
\item 对于上述肥皂泡，使其带上电量 $Q$，如果肥皂泡带电后重新达到平衡的过程为等温过程，达到平衡后泡的半径为 $R$，泡内气体可以近似为理想气体，试确定现在泡内气体的压强，并从物理上定性说明 $R$ 将变大还是变小。
\item 试给出 $R$ 以 $p_0,\sigma,Q$ 和 $R_0$ 表示的形式，并比较 $R$ 和 $R_0$ 的大小。
\end{enumerate}}
\answerarea
\problem{6.16}{对在空气中匀速下落的雨滴测得其上下两端的间距为 $d=5.0\,\mathrm{mm}$，试确定该雨滴上下两端处的曲率半径之差。}
\answerarea
\clearpage
\problem{6.17}{水深为 $h=2.0\,\mathrm m$ 的池底产生许多直径为 $d_0=5.0\times10^{-5}\,\mathrm m$ 的小气泡，如果水的表面张力系数 $\sigma=7.3\times10^{-2}\,\mathrm{N/m}$，大气压强 $p_0=1.01\times10^5\,\mathrm{Pa}$，试确定这些气泡等温地上升到水面时，它们的直径各为多大？}
\answerarea
\problem{6.18}{图示是测定表面张力系数的一种装置，慢慢提起放入待测液体中的铜片，在它刚要离开液面时，测得所用的上提力 $F$，即可测出表面张力系数。设待测液体与铜片的接触角 $\theta=0$，铜片质量 $m=5.0\times10^{-4}\,\mathrm{kg}$，宽度 $l=3.977\times10^{-2}\,\mathrm m$，厚度 $d=2.3\times10^{-4}\,\mathrm m$，$F=1.07\times10^{-2}\,\mathrm N$，该实验测出的液体的表面张力系数为多大？
\fig{6_18}}
\answerarea
\clearpage
\problem{6.19}{如图所示的 U 形管的两侧细管的口径不同，较细的管 $A$ 的内半径 $r=5.0\times10^{-5}\,\mathrm m$，较粗的管 $B$ 的内半径 $R=2.0\times10^{-4}\,\mathrm m$。向该 U 形管中注入适量的表面张力系数为 $\sigma=7.3\times10^{-2}\,\mathrm{N/m}$ 的水，试确定两侧管中水面的高度差 $h$。
\fig{6_19}}
\answerarea
\problem{6.20}{将少量水银放在两块水平的平玻璃板之间，已知水银的表面张力系数 $\sigma=0.45\,\mathrm{N/m}$，水银与玻璃的接触角 $\theta=135^\circ$，试确定多大的负荷加在上板时，能使两板间的水银厚度处处都等于 $1.0\times10^{-4}\,\mathrm m$，并且两板和水银的接触面积都为 $4.0\times10^{-3}\,\mathrm{m^2}$。}
\answerarea
\clearpage
\problem{6.21}{内半径为 $R_1=2.0\times10^{-2}\,\mathrm m$ 的玻璃管中，插入一半径为 $R_2=1.5\times10^{-3}\,\mathrm m$ 的同轴玻璃棒。已知水的密度 $\rho=1.00\times10^3\,\mathrm{kg/m^3}$，表面张力系数 $\sigma=7.3\times10^{-2}\,\mathrm{N/m}$，水与玻璃的接触角 $\theta=0^\circ$，试确定水在管中上升的高度。}
\answerarea
\problem{6.22}{玻璃管的内直径 $d=2.0\times10^{-5}\,\mathrm m$，长为 $l=0.20\,\mathrm m$，垂直插入水中，管的上端是封闭的。已知大气压强 $p_0=1.013\times10^5\,\mathrm{Pa}$，水的表面张力系数 $\sigma=7.3\times10^{-2}\,\mathrm{N/m}$，水与玻璃的接触角 $\theta=0^\circ$，试问插入水面下的那一段长度应为多少，才能使管内外水面一样高？}
\answerarea
\clearpage
\problem{6.23}{将一充满水银的气压计下端浸在一个广阔的盛有水银的容器中，读数为 $p=0.950\times10^5\,\mathrm{Pa}$。已知毛细管的直径 $d=2.0\times10^{-3}\,\mathrm m$，接触角 $\theta=\pi$，水银的表面张力系数 $\sigma=0.49\,\mathrm{N/m}$。
\begin{enumerate}
\item 试确定水银柱的高度；
\item 考虑到毛细现象后，真正的大气压强为多少？
\item 如果允许误差为 $0.1\%$，毛细管直径所能允许的最小值为多大？
\end{enumerate}}
\answerarea
\problem{6.24}{内径 $d=4.0\times10^{-4}\,\mathrm m$、长 $l_0=0.20\,\mathrm m$ 的玻璃管，水平且等温地浸到深度 $h=0.15\,\mathrm m$ 的水银槽中，空气全部留在管中。已知大气压强 $p_0=1.013\times10^5\,\mathrm{Pa}$，水银的表面张力系数 $\sigma=0.490\,\mathrm{N/m}$，水银与玻璃之间的接触角 $\theta=\pi$，试确定玻璃管中空气柱的长度 $l$。}
\answerarea
\clearpage
\problem{6.25}{一根内直径为 $1\,\mathrm{mm}$ 的玻璃管，竖直插入盛水银的容器里，管的下端在水银面以下 $1\,\mathrm{cm}$ 处。
\begin{enumerate}
\item 要在管的下端吹出一半球形气泡，管中气体的计示压强是多少？所谓计示压强是指实际压强与大气压 $p_0$ 之差；
\item 如果管内气压比大气压低 $3000\,\mathrm{N/m^2}$，水银与玻璃的接触角是 $135^\circ$，水银在管内会升到多高？
\end{enumerate}}
\answerarea
\problem{6.26}{一半径为 $r$、密度为 $\rho$ 的均质扁圆柱体，浮在密度为 $\rho_1$ 的某种液体中，浸入液体部分的高度为 $h$，高出液面的高度为 $H$，纵切面如图所示。若已知液体与固体间的表面张力系数为 $\sigma$。
\begin{enumerate}
\item 试就液体与固体间的接触角 $\theta=0$（完全润湿）和 $\theta=\pi$（完全不润湿）两种情况求出高度 $h$。
\item 如果表面张力系数随温度按规律 $\sigma=\sigma_0(1-kt)$ 变化，其中 $\sigma_0$ 和 $k$ 是常量，$t$ 是摄氏温标下的温度，试确定，当温度变化时，为保持浸入深度 $h$ 不变，液体的体膨胀系数随温度应怎么变化？
\end{enumerate}
\fig{6_26}}
\answerarea
\clearpage
\problem{6.27}{在半径 $r=0.3\,\mathrm{mm}$ 的毛细管中注水，结果在管下端形成一个水滴，如图所示。已知表面张力系数 $\sigma=7.12\times10^{-2}\,\mathrm{N/m}$，若把水滴表面看成一个半径为 $R=3\,\mathrm{mm}$ 的球面的一部分，则水柱的高度 $h$ 为多少？
\fig{6_27}}
\answerarea
\problem{6.28}{有一开口向上、竖直放置的 U 形玻璃毛细管，内半径为 $r=1\,\mathrm{mm}$，水平连通部分长为 $5\,\mathrm{cm}$，现在 U 形管中注入 $2\,\mathrm{cm}$ 长的水银柱及 $30\,\mathrm{cm}$ 长的水柱。已知水的表面张力系数为 $0.07\,\mathrm{N/m}$，它与玻璃完全润湿；而水银的表面张力系数为 $0.5\,\mathrm{N/m}$，密度为 $13.6\times10^3\,\mathrm{kg/m^3}$，它与玻璃完全不润湿。如果不考虑水银与水接触面处液面的弯曲，在 U 形管中液面的高度差是多少？}
\answerarea
\clearpage
\problem{6.29}{在一根两端开口的毛细管中滴入一滴水后将它竖直放置，若滴入的水在毛细管中分别形成长为（1）$2.00\,\mathrm{cm}$，（2）$4.00\,\mathrm{cm}$，（3）$2.98\,\mathrm{cm}$ 的水柱，毛细管的内（直）径为 $1.0\,\mathrm{mm}$，并且能够完全被水润湿，水的表面张力系数为 $0.073\,\mathrm{N/m}$，试确定在上述三种情况下，水柱的上、下液面各呈什么样的形状。}
\answerarea
\problem{6.30}{试确定当可能的毛细现象引起的上升力与重力达到平衡时，液体表面的曲线方程（可仅考虑两维情况）。}
\answerarea
\clearpage
\problem{6.31}{两块平行且竖直放置的玻璃板部分浸入水中，如果板的长度为 $15\,\mathrm{cm}$，两板之间的间距为 $0.1\,\mathrm{mm}$，水的表面张力系数为 $0.07\,\mathrm{N/m}$，试确定在水能够完全润湿玻璃板的情况下，两板之间的吸引力。}
\answerarea
\problem{6.32}{一无限大平板竖直地部分浸入液体中，记液体的密度为 $\rho$、表面张力系数为 $\sigma$，假设液体完全润湿平板，试确定液体在此平板板面上上升的高度。}
\answerarea
\clearpage
\problem{6.33}{如果农作物、花草灌木及树木中传输水分及营养物质的导管为内半径为 $5.0\,\mu\mathrm m$ 的细圆柱形管，水分及营养物质的表面张力系数为 $9.0\times10^{-2}\,\mathrm{N/m}$、密度为 $1.1\times10^3\,\mathrm{kg/m^3}$，地面及其附近的重力加速度为 $9.8\,\mathrm{m/s^2}$，试通过定量计算说明水分及营养物质在农作物及花草灌木中传输的机制可能主要是毛细现象，而在高大的树木中一定还有其它机制。}
\answerarea
\clearpage
